% ======================================================================
% Section 6 — The pq side: the tensor translation
% ======================================================================

\section{The $pq$ side: the tensor translation}\label{sec:pq}

The group-ring translation of \S\ref{sec:groupring} was stated and verified at
prime modulus, where the three covering objects are APs of $\Z_p$. This
section completes the structural picture at semiprime modulus: at $N=pq$ the
CRT identification $\Z_{pq}\cong\Z_p\times\Z_q$ turns the two-unit and
three-unit cells into \emph{fiber-wise} three-AP problems of exactly the
prime-modulus form, with one member fixed by the kernel side and the other
members moving affinely along the fiber index. The tensor structure that the
outlook of \S\ref{sec:gt-outlook} anticipated is real, it is proved
(Theorem~\ref{thm:R}), and it is machine-verified in every component. The
outcome is the structural completion the programme wanted: \emph{the $pq$
cells introduce no new computational input} --- they close modulo exactly the
two lemmas already named for the $p^2$ closure (Corollary~\ref{cor:pqdep}).

\subsection{Setup: the CRT picture and the closed-form constants}
\label{sec:pq-setup}

Fix $N=pq$ with $5\le p<q$ primes, $c=\floor{(N-1)/6}$, and work in the CRT
coordinates $(x,y)=(k\bmod p,\ k\bmod q)$. Define
\[
\rho_p:=\floor{c/q}=\floor{p/6},\qquad \rho_q:=\floor{c/p}=\floor{q/6},
\]
the first equality by the nested-floor identity. A kernel-$p$ member (speed
$pa$, $a\in\Z_q^{*}$) has $\|pa\,k\|_{N}=p\,\|ak\bmod q\|_q\le c$ iff
$\|ak\bmod q\|_q\le\rho_q$: it covers exactly the rows of
\[
Y_a=\{y\in\Z_q:\|ay\|_q\le\rho_q\},
\]
an AP of $2\rho_q+1$ rows (the preimage of the symmetric interval under
multiplication by $a$); symmetrically a kernel-$q$ member (speed $qb$,
$b\in\Z_p^{*}$) covers exactly the columns of
$X_b=\{x\in\Z_p:\|bx\|_p\le\rho_p\}$, an AP of $2\rho_p+1$ columns. The
uncovered region of a $(1,1,2)$ family --- one kernel of each type, two units
--- is the rectangle $X_b^{c}\times Y_a^{c}$ with
\[
|X_b^{c}|=p-2\rho_p-1,\qquad |Y_a^{c}|=q-2\rho_q-1\ \tfrac{2}{3}q-\mathrm{O}(1).
\]

A unit $u\in\Z_{pq}^{*}$ has bad set $B_u=u^{-1}[-c,c]$, and its trace on a
row is governed by one integer interval. Write $k=y+qj$ ($j\in\Z_p$) for the
row $y$ and $\sigma=\sigma_u(y):=uy\bmod q\in[0,q)$; then
$(x,y)\in B_u$ iff $x\equiv u^{-1}t\pmod p$ for some $t\in[-c,c]$ with
$t\equiv\sigma\pmod q$, i.e.\ $t=\sigma+qj'$ with
\begin{equation}\label{eq:Jinterval}
j'\in J_u(y):=\bigl[\lceil(-c-\sigma)/q\rceil,\ \floor{(c-\sigma)/q}\bigr],
\end{equation}
an interval of $n_u(y):=|J_u(y)|=\floor{(c-\sigma)/q}+\floor{(c+\sigma)/q}+1$
consecutive integers. This is the $pq$ form of the fiber footprint
$T_u(j)$ of the $p^2$ analysis.

\subsection{The row translation}\label{sec:pq-thm}

\begin{theorem}[row translation at $N=pq$]\label{thm:R}
With the notation above, writing $\alpha_u:=u^{-1}\bmod p$:
\begin{enumerate}[leftmargin=1.9em,itemsep=2pt]
\item[\textup{(i)}] the trace of $B_u$ on the row $y$ is the AP
\[
f_u(y)=\{\alpha_u(\sigma+qj')\bmod p:\ j'\in J_u(y)\},
\]
of step $\alpha_u q\bmod p$ and size $n_u(y)$;
\item[\textup{(ii)}] the size profile is class-appropriate: writing
$p=6k+\epsilon$ with $\epsilon\in\{1,5\}$,
\[
n_u(y)\in
\begin{cases}
\{2k+1,\ 2k+2\}&\epsilon=5,\\
\{2k,\ 2k+1\}&\epsilon=1,
\end{cases}
\]
independently of $q$;
\item[\textup{(iii)}] the $(1,1,2)$ family covers $\Z_{pq}$ if and only if,
for every row $y\in Y_a^{c}$,
\[
X_b^{c}\ \subseteq\ f_{u_1}(y)\cup f_{u_2}(y),
\]
i.e.\ iff $(f_{u_1}(y),\,f_{u_2}(y),\,X_b)$ is a three-AP covering of
$\Z_p$ whose third member is the \emph{fixed} AP $X_b$; at $\epsilon=5$ this
is the exact partition of the multiset $\{2k{+}2,2k{+}2,2k{+}1\}$, and at
$\epsilon=1$ a critical covering with $\mathrm{ov}\in\{0,1,2\}$ and sizes in
$\{2k,2k{+}1\}$ --- in both cases precisely the census scope of
\S\ref{sec:gt-census} with $s_1=|X_b|=2k+1$;
\item[\textup{(iv)}] the start $a_u(y)=\alpha_u(\sigma_u(y)+q\,j_{\min}(y))$
of the row AP moves \emph{piecewise-affinely} in $y$:
$a_u(y{+}1)-a_u(y)\in\{\alpha_u u,\ \alpha_u(u-q)\}$ in $\Z_p$; the
$(1K,3U)$ cells obey the same translation with three footprints and no fixed
member (the whole row must be covered), i.e.\ the Lemma~A/B problem itself.
\end{enumerate}
\end{theorem}

\begin{proof}
(i) is the computation above \eqref{eq:Jinterval}. (ii) Write
$c=\rho_p q+c'$ with $c'=c\bmod q$. Since $p=6m+\epsilon$ and
$r:=(\epsilon q)\bmod 6\in\{1,5\}$,
$c=mq+(\epsilon q-r)/6$, so $c'=(\epsilon q-r)/6$, which is $>q/2$ exactly
when $\epsilon=5$ (for $q\ge 3$) and $<q/2$ exactly when $\epsilon=1$. Now
\[
n_u(y)=2\rho_p+\floor{(c'-\sigma)/q}+\floor{(c'+\sigma)/q}+1,
\]
where the two floors take values in $\{-1,0\}$ and $\{0,1\}$: the maximum
$2\rho_p+2$ occurs iff $q-c'\le\sigma\le c'$ (nonempty iff $c'>q/2$, i.e.\
$\epsilon=5$), the minimum $2\rho_p$ occurs iff
$c'<\sigma<q-c'$ (nonempty iff $c'<q/2$, i.e.\ $\epsilon=1$), and
$2\rho_p+1$ always occurs. With $\rho_p=\floor{p/6}=k$ this is the stated
profile; the dichotomy depends on $c'$, hence on $p\bmod 6$ and $q$, but the
\emph{window} $\{2k{+}1,2k{+}2\}$ resp.\ $\{2k,2k{+}1\}$ depends only on
$p$. (iii) The uncovered region is exactly $X_b^c\times Y_a^c$; covering it
is equivalent to the row inclusions. At $\epsilon=5$,
$|X_b^{c}|=p-2k-1=4k+4=2(2k+2)$ while $n_{u_i}(y)\le 2k+2$, so both
footprints must have the maximal size and
$f_{u_1}(y)\sqcup f_{u_2}(y)=X_b^{c}$ exactly: together with $X_b$ (disjoint
from $X_b^{c}$) this is a pairwise-disjoint partition of $\Z_p$ with sizes
$(2k+2,2k+2,2k+1)$, the B-side critical multiset with the third member fixed.
At $\epsilon=1$, $|X_b^{c}|=4k$ and $n_{u_i}(y)\le 2k+1$, so
$n_1+n_2+|X_b|-p\le 2$ with all sizes in the A-side window: a critical
covering with $s_1=2k+1$ fixed. (iv) $\sigma_u(y{+}1)-\sigma_u(y)\in\{u
\bmod q,\ u\bmod q-q\}$ and $j_{\min}(y{+}1)-j_{\min}(y)\in\{0,-1\}$, giving
the two rates; the $(1K,3U)$ statement is the same row decomposition with the
kernel covering only rows.
\end{proof}

Part (iii) is the tensor translation: the $pq$ cell condition \emph{is} the
prime-modulus census condition, fiber by fiber, with the kernel side
supplying the fixed member and the unit side supplying the moving members.
Part (ii) is what makes it sharp --- the row-size profile is forced by
$c\bmod q$ to sit in the class-appropriate critical window, so no
out-of-scope size ever appears on a fiber. Part (iv) is the shear structure:
the covering configurations of the rows are pinned (below), and their
carriers drift at rates determined by the units alone.

\subsection{The row census and the pinning}\label{sec:pq-census}

Theorem~\ref{thm:R}(iii) reduces the row condition to the census of
\S\ref{sec:gt-census} \emph{restricted to} $s_1=|X_b|=2k+1$. That restricted
census is again four families with $p$-independent counts.

\begin{proposition}[row census, machine]\label{prop:rowcensus}
At every prime $p\equiv5\pmod6$ with $11\le p\le59$, the census of critical
coverings in the normal form with $s_1=2k+1$ (sizes $(2k{+}2,2k{+}2)$,
$\mathrm{ov}=0$) returns exactly $24$ configurations, distributed by the
$\pm$-rep pair $(\|d_2\|,\|d_3\|)$ as
\[
(1,1):8,\qquad(2,2):8,\qquad(1,q):4,\qquad(q,1):4,
\]
with $q=(p-1)/2$. At every prime $p\equiv1\pmod6$ with $19\le p\le61$, the
census with $s_1=2k+1$ and sizes in $\{2k,2k+1\}$ ($\mathrm{ov}\in\{0,1,2\}$)
returns exactly $312$, distributed as
\[
(1,1):104,\qquad(2,2):72,\qquad(1,q):68,\qquad(q,1):68.
\]
All counts are independent of $p$, and no other family occurs.
\end{proposition}

The $24$ of the $\epsilon=5$ side is completely proved by the parametrization
mechanism of Proposition~\ref{prop:pindep}: the $(1,1)$ configurations are the two
cyclic orders of the consecutive tiling of $I_1=[2k+1,p-1]$ by two intervals
of $2k+2$, in $2^2$ parametrizations ($8$); the $(2,2)$ configurations are
the even/odd split of $I_1$ in both assignments and orientations ($8$); and
the $(1,q)$, $(q,1)$ configurations are the \emph{wrapped} two-block: the
two-block of step $-q$ with first block $[2k+1,3k+1]$ has its second block at
$[a-q,a-q+k]\equiv[5k+4,6k+4]$ through the wrap at $0$, and its complement
in $I_1$ is the single interval $[3k+2,5k+3]$ of size $2k+2$ --- one
set-arrangement in $2^2$ parametrizations ($4$ each). The unwrapped two-block
placements $a\in\{3k+1,3k+2\}$ fail by exactly one exposed endpoint
($\{2k+1\}\cup[3k{+}3,5k{+}3]$ resp.\ $[3k{+}2,5k{+}2]\cup\{6k{+}4\}$): the
near-miss that the shear argument exploits. This is the CRT-compatible form
of the four families: on each fiber, the classified objects are exactly the
census objects, with the fixed member playing the normalized role.

The pinning statement is the machine result that connects this to
Input~S. Across the eight moduli
$N\in\{35,55,77,91,95,115,119,143\}$, all $132{,}485$ $(a,b,u_1,u_2)$
families were scanned row by row:

\begin{proposition}[pinning at $N=pq$, machine]\label{prop:pqpinned}
Every row configuration $(f_{u_1}(y),f_{u_2}(y),X_b)$ satisfying the row
condition of Theorem~\emph{\ref{thm:R}(iii)} --- $37{,}672$ of them across
the eight moduli --- is, after the normalization $x\mapsto bx+\rho_p$,
\emph{exactly} a census solution of Proposition~\ref{prop:rowcensus} at the
row modulus ($100\%$; both the family label and the placement sets coincide).
For a fixed family $(a,b,u_1,u_2)$ the number of rows satisfying the
condition is at most $6$.
\end{proposition}

This is the identity-level form of the shear mechanism at $N=pq$: a covering
would need the row condition on \emph{every} row of the arc $Y_a^{c}$ (of
size $q-2\floor{q/6}-1\ \tfrac23 q\ \ge 8$ for $q\ge11$), but the pinned
placement sets have $p$-independent size ($24$ resp.\ $312$ elements, by
Propositions~\ref{prop:rowcensus} and~\ref{prop:pindep}), and the carriers
drift at the two rates of Theorem~\ref{thm:R}(iv). Either the relative drift
is nonzero, and the normalized placement pair exits the pinned set within a
bounded number of rows (at most $6$ in every measured case, with the failure
rows carrying $1$--$2$ uncovered points each --- the margins below), or the
relative drift vanishes, which forces a multiplicative relation between
$u_1,u_2$ modulo $p$ --- the collapse case, where the family degenerates to a
common scaling and the kernel side must cover. The dichotomy is the same as
at $p^2$, with the same two consumers.

\subsection{Machine verification and the cell record}\label{sec:pq-verify}

\begin{table}[t]
\centering
\small
\caption{Verification of the $pq$ translation and the cells. ``Families''
counts $(a,b,u_1,u_2)$ resp.\ $(a,u_1,u_2,u_3)$ tuples over $\pm$-canonical
speeds; margins are the minimal numbers of uncovered points over all
families.}
\label{tab:pqverify}
\begin{tabularx}{\textwidth}{@{}lXXr@{}}
\toprule
check & scope & outcome & count \\
\midrule
constants & $8$ moduli $35$--$143$ & $\rho_\cdot=\floor{\cdot/6}$;
$\epsilon{=}5$ partition equality $|X^c|=2(2k{+}2)$; profile dichotomy
\eqref{eq:Jinterval}--(ii) & 8 \\
footprint formula & all units $\times$ all rows, $8$ moduli & exact set
equality $f_u(y)=\{\alpha_u(\sigma+qj')\}$ & 4{,}394 \\
row census (Prop.~\ref{prop:rowcensus}) & $\epsilon=5$: $p\le59$;
$\epsilon=1$: $p\le61$ & $24$ resp.\ $312$, counts constant, four families
only & 14 \\
$(1,1,2)$ cell & $N\in\{35,\dots,143\}$, all $(a,b,u_1<u_2)$ & $0$
coverings; margins $4$--$18$ & 143{,}733 \\
$(1K,3U)$ cells, both orientations & same $N$, all unit triples & $0$
coverings; margins $2$--$10$ & 989{,}000 \\
pinning (Prop.~\ref{prop:pqpinned}) & $8$ moduli, all families & $37{,}672/
37{,}672$ rows are census solutions; max good rows $\le6$ & --- \\
$T$-10 reproduction & $N=77,91$ & $(1,1,2)$ empty, as committed & 2 \\
parity/two-block structure & every census $(2,2)$ solution, $6$ primes &
one integer parity class of $I_1$, $\le\mathrm{ov}$ spills & 1{,}248 \\
\bottomrule
\end{tabularx}
\end{table}

All checks pass (Table~\ref{tab:pqverify}); the harness is an independent
code path (formula-based footprints cross-checked against the direct
definition, bit-masked cell enumeration, census by residual fit with the
$p=7$ brute-force validation of \S\ref{sec:gt-verify} carried over). The
margin structure mirrors the prime side: the nearest $(1,1,2)$ families sit
in the interval and parity families with a handful of good rows (the drift
kills the rest, each failing row missing $1$--$2$ points of $X_b^c$), and the
nearest $(1K,3U)$ families fail entire rows by $2$--$10$ points --- the
cyclotomic analogue of the shear near-miss at $p^2$ (families failing all
but two fibers).

\begin{corollary}[the $pq$ closure's dependencies]\label{cor:pqdep}
The emptiness of the $(1,1,2)$ and $(1K,3U)$ cells at $N=pq$ ($5\le p<q$),
verified here for $N\le143$, is modulo exactly the computational inputs of
the $p^2$ closure, and no new unbounded statement appears on the $pq$
side: \emph{(a)} the support lemma (Conjecture~\ref{conj:SL}) at the row
and column moduli --- which classifies the per-fiber configurations into
the four families with the pinned placement sets of
Propositions~\ref{prop:rowcensus}--\ref{prop:pqpinned} --- and \emph{(b)}
the shear-rigidity argument applied along the row arc
(Theorem~\emph{\ref{thm:R}(iv)}), i.e.\ Input~S in its fiberwise form. The
$p^2$ closure's foot-count input (Input E) reaches the $pq$ side only
through the prime-side trichotomy inside \emph{(a)} and is now
discharged (Lemma~\ref{lem:E}); after this revision, both closures rest
on the single unbounded input S.
\end{corollary}

\begin{proof}[Proof sketch]
By Theorem~\ref{thm:R}(iii) a covering family yields a critical three-AP
covering on every row of $Y_a^{c}$; by the support lemma each is one of the
four families in the normalization of the fixed member $X_b$, with placement
in the pinned set; the affine drift (iv) moves the normalized placements at
rates $\lambda_i\in\{\alpha_{u_i}u_i,\alpha_{u_i}(u_i-q)\}$, and the pinned
set is finite and $p$-independent while $|Y_a^{c}|$ grows with $q$; a
nonvanishing relative drift exits the set (Proposition~\ref{prop:pqpinned}
measures at most $6$ surviving rows), and a vanishing one collapses the units
to a common scaling modulo the fiber modulus, whereupon the reduced problem
is the prime-modulus cell, closed by Lemmas A/B. The column side exchanges
the roles of $p$ and $q$ and of rows and columns.
\end{proof}

The reformulation is thus structurally complete on the semiprime side: every
cell of the two-unit and three-unit counting at $N=pq$ is either closed by
the proved fiber-cap theorems (Theorems 2--3 of the companion notes, with the
$(1,1,2)$ cell's row-slack equality $p-2\rho_p-1=2(2\rho_p+2)$ at
$p\equiv5\pmod6$ explaining precisely why that cell is the sole survivor of
the counting), or reduced by Theorem~\ref{thm:R} to the prime-modulus
programme with the same named inputs --- of which the foot-count lemma is
now proved (Lemma~\ref{lem:E}), leaving a single unbounded input. What
remains outside the
classification is the rigid-zone lift question, unchanged.
